
Provenance is not the only place public-key cryptography touches media.
Archives of medical imaging, evidence and surveillance footage encrypt each
object under its own data key and store that key wrapped under a public key,
and harvest-now-decrypt-later applies to them without qualification. Migrating
such an archive means unwrapping every data key under the legacy key and
wrapping it again. For the post-quantum wrap we used the shape of the
construction RFC~9629 specifies for CMS~\cite{rfc9629}: ML-KEM gives a shared
secret, HKDF-SHA256 derives a key-encryption key from it, and AES key
wrap~\cite{rfc3394} wraps the 32-byte data key into 40 bytes. We did not build
the full \texttt{KEMRecipientInfo}: the key derivation uses a fixed
\texttt{info} string rather than the DER-encoded
\texttt{CMSORIforKEMOtherInfo} the RFC requires, and no recipient structure is
serialised, so the measurement shows the cost of the cryptography and not
conformance to the RFC.

Re-wrapping protects the archive from the day it finishes. It does nothing
for copies of the legacy wrapped keys an adversary has already harvested,
which a future quantum computer will open whatever the archive does
afterwards; for data that must stay confidential for decades, the date that
matters is the one the wrapped keys were first exposed.

Table~\ref{tab:keywrap} gives the cost per key. Storage grows by a factor of
4.4, from 256 bytes under RSA-2048 to 1,128 under ML-KEM-768, which is the
1,088-byte KEM ciphertext and the 40-byte wrapped key; a CMS container adds a
few dozen bytes of its own fields to either. Time runs the
other way. An RSA-2048 unwrap takes 0.53\,ms and a complete ML-KEM-768 wrap
0.017\,ms, of which key derivation and key wrap are 0.003\,ms, so the legacy
unwrap is 97\% of the cryptographic work in a re-wrap. Once an archive has
migrated, unwrapping is 24 times faster than it was; as Section~\ref{sec:hash} of the paper found for signing, the post-quantum operation is not the expensive part.


Cost per key does not give the cost of a migration, because the I/O depends on
where the wrapped key is kept. Table~\ref{tab:rewrap} re-wraps archives on
disk in three layouts. When the key sits in a sidecar file, or in a header
slot that was sized for the post-quantum key when the object was written, the
migration runs at 1,000 to 1,500 objects per second and the RSA unwrap takes
between 63\% and 82\% of each object's time, whatever the object's size. When
the key sits in a header sized for RSA, the larger key cannot be written in
place and the whole object is rewritten. For photographs of about 100\,KiB
that adds about 30\% per object against the reserved layout. For 4\,MiB
objects re-wrapping is 4.2 times slower, and reading and rewriting the object
take 77\% of the time. The penalty grows with the object, and archives of
scans and footage hold objects far larger than 4\,MiB. The lesson is the one
C2PA's reserved signature block already teaches
(Section~\ref{sec:method-decisions} of the paper): a format that leaves room for a larger
key migrates in place, and one that does not rewrites its content.


Those reads came from a warm cache, since each archive was written just
before it was re-wrapped, and the writes were forced with \texttt{fsync},
which on macOS does not flush the drive's own cache. A second run emptied the
page cache first and forced every write with \texttt{F\_FULLFSYNC}, which does
(Table~\ref{tab:rewrapcold}). Throughput falls to 101 to 167 objects a
second, because the flush now takes 4.4 to 7.0\,ms of each object's 5.2 to
9.1, and the order of the layouts holds: the reserved slot is the fastest at
both sizes, and rewriting a 4\,MiB object takes 1.7 times as long as updating
it in place. The RSA unwrap also slows, to 0.6 to 1.9\,ms, which we attribute
to the processor clocking down while it waits on each flush. One limit
remains: an in-place update of a reserved slot is not atomic, and a production
format would write the new key to a second slot and switch to it, which we did
not measure; the in-place reserved-slot figures should be read as a lower
bound for a crash-safe layout.

